\section{The determinant-fixed Quot obstruction theory}\label{sec:derived}

We make the virtual construction in two settings:
\begin{enumerate}[label=(\roman*)]
\item $X$ is a smooth connected projective complex surface with $p_g=q=0$;
\item $X$ is a smooth connected projective complex surface with $p_g=q=1$,
      $K=K_X$ ample, and $D=3K$.
\end{enumerate}
Write $K=K_X$.  In both settings $\chi(\mathcal O_X)=1$.
The second setting is used only to determine universal coefficients in
Section~\ref{sec:rankone}.  No positivity assumption on $K$
is imposed in setting~(i).  All virtual classes and equivariant Chow groups
in this Section have rational coefficients.

Our goal is an equivariant perfect obstruction theory on the
determinant-fixed Quot scheme, together with its actual restriction to
the fixed locus.  Surface Quot obstruction theories and localization
relations for a trivial target bundle were developed by
Marian, Oprea, and Pandharipande \cite[Section~1]{MOP}; subsequent
intersection-theoretic and $K$-theoretic developments appear in
\cite{OP,AJLOP}.  Here the based determinant fiber must also work on
the irregular auxiliary surface.  We construct that fiber as a derived
moduli problem so that its obstruction morphism, and not only its
virtual tangent class, is specified.

Fix line bundles $L_1,L_2,D$, and put $V=L_1\oplus L_2$ and $A=L_1-L_2$.
For an integer $n$, let $B$ be the scheme of quotients
\begin{equation}\label{der:quot-sequence}
\begin{gathered}
0\longrightarrow S'\xrightarrow{i}V\longrightarrow Q\longrightarrow0,
\\
\rk Q=0,\quad \det Q=D,\quad
\int_X\ch_2(Q)=n+L_1D-\frac{D^2}{2}.
\end{gathered}
\end{equation}
Here the determinant condition is equality in the Picard scheme.
Choose a polarization on $X$.  The stated Chern data determine a Hilbert
polynomial, and $B$ is the determinant fiber in the corresponding projective
Quot scheme.  In particular, $B$ is projective.  The group
$G=\mathbb G_m$ acts on $V$ with characters $+1,-1$ on its two summands,
and hence acts on $B$.  We write $t=c_1^G(\mathbb C_1)$ for the equivariant
parameter.  Later we specialize $t$ to $1/2$, so the evaluated weights of
the two summands are $w_1=1/2$ and $w_2=-1/2$.

Our complexes use cohomological grading: a sheaf $F[-1]$ lies in degree
one.  Pullbacks and tensor products of complexes are derived unless
specified otherwise.  If $Y$ is a parameter scheme and
$\pi:X\times Y\to Y$ is projection,
we abbreviate
\[
\RHom(F,G)=R\pi_*R\mathcal{H}om(F,G).
\]
Thus Hom complexes in families are complexes on the parameter scheme.
The notation $\RGamma(F)$ for an object on $X$ denotes its derived sections,
pulled back as a constant complex when parameters are present.

\subsection{Derived arrows and classical families}

Ciocan-Fontanine and Kapranov constructed a derived Quot scheme
whose truncation is the classical Quot scheme and whose tangent
cohomology is given by Ext groups from the kernel to the quotient
\cite[Theorem~4.3.2]{CFK}. Here we use a perfect-arrow moduli
description to retain the maps needed for the based determinant fiber.

We use the derived moduli stack of perfect complexes and its morphism
stack as constructed in \cite[Theorem~3.6, Definition~3.18, and Corollaries~3.17,
3.19, and~3.29]{TV}.  For smooth proper $X$ these stacks are locally
geometric and locally of finite presentation.  The morphism stack with
target fixed to $V$ also has a direct description: over a perfect source
$F$ its fiber is the derived linear mapping stack associated to
$\RHom(F,V)$.  This complex is perfect, and its formation commutes with
derived base change, because $X$ is smooth and proper.

\begin{lemma}[The open Quot enhancement]\label{der:arrow}
The fixed-target derived arrow stack has an open substack $\mathfrak U$
whose truncation $U$ is the ordinary Quot scheme with the chosen rank and
Hilbert polynomial, before fixing determinant.  On $U$, its tangent
complex is naturally
\begin{equation}\label{der:unrestricted-tangent}
\begin{aligned}
T_{\mathfrak U}|_U
&\simeq\Cone\bigl(\RHom(S',S')\xrightarrow{i\circ-}\RHom(S',V)\bigr)\\
&\simeq\RHom(S',Q).
\end{aligned}
\end{equation}
It is perfect of Tor-amplitude $[0,1]$.  These identifications are
compatible with forgetting the source or target, with direct sums of
arrows, and with base change.  The same assertions hold for a fixed
vector-bundle target of any rank and torsion quotients.
\end{lemma}

\begin{proof}
First take the locus where both the source and its cone are sheaves in
degree zero on classical fibers, with the prescribed Hilbert polynomials.
We explain the family conditions in this description.  For a perfect
complex on $X\times Y$, the locus where it has relative Tor-amplitude zero
over $Y$ is open.  Locally one represents the complex by vector bundles.
At a fiber on which it is exact away from degree zero, the positive end
can be split off.  At the negative end, fiberwise injectivity between
modules flat over the base implies injectivity with base-flat cokernel,
by the local criterion of flatness.  Successively replacing the negative
end by these cokernels gives a sheaf flat over the base near that fiber.
Openness of flatness gives an open neighborhood in $X\times Y$; properness
of $X\times Y\to Y$ then gives the required open neighborhood in $Y$.
Applying this to the source and cone proves the asserted openness.

Conversely, let $V_Y\twoheadrightarrow Q$ be a flat quotient family over
a noetherian base $Y$.  Its kernel $S'$ is flat over $Y$.  Both $Q$ and
$S'$ are perfect on $X\times Y$, even if $Y$ is singular.  Indeed, resolve
either sheaf locally by finite-rank free modules.  Every successive
syzygy remains flat over $Y$.  On a surface fiber the local rings are
regular of dimension at most two, so after two steps the syzygy is locally
free on that fiber.  A finitely presented, base-flat module with locally
free fibers is locally free near the point, again by the local flatness
criterion.  This truncates the resolution.  This argument suffices over
noetherian bases, and finite presentation extends the moduli
identification to arbitrary classical test schemes.
Flat approximation and the perfection criteria in
\cite[Tags 02JO, 0DHV, and 068X]{Stacks} also give this last passage:
a finitely presented base-flat module is pseudo-coherent, and the
regular surface fibers give absolute Tor-amplitude contained in $[-2,0]$.

On our open locus the cone sequence is consequently an exact sequence
of sheaves with flat quotient.  An automorphism $a$ of its source
satisfying $i\circ a=i$ is the identity, since $i$ is injective.
Negative Ext groups between these sheaves vanish as well, so no higher
automorphisms remain in the classical moduli problem.  The truncation
is therefore the Quot scheme, including its functor on nonreduced bases.

Here is a tangent calculation that also records the relevant maps.
Allow both vertices of an arrow $F\to G$ to vary.  Its derived
endomorphism complex is
\[
\fib\left\{
\RHom(F,F)\oplus\RHom(G,G)
\xrightarrow{(a,b)\mapsto f a-b f}\RHom(F,G)
\right\}.
\]
One can compute this complex with the two-term resolution of a one-arrow
diagram by diagrams induced from its vertices.  More explicitly, if
$P_1=(\mathcal O_X\xrightarrow{\id}\mathcal O_X)$ and
$P_2=(0\to\mathcal O_X)$, the resolution has terms
$P_2\otimes F$ and $(P_1\otimes F)\oplus(P_2\otimes G)$.
Its map at the target vertex sends $s$ to $(s,-f(s))$.
Thus perfect vertices give a perfect diagram.  The one-arrow path algebra
is finite dimensional and has this finite projective resolution; the
usual perfect-object tangent formula applies to its diagrams on $X$.
The tangent of their moduli is the displayed endomorphism complex
shifted by $[1]$.

Fixing the target takes the tangent fiber over
$\RHom(V,V)[1]$.  The result is the first cone in
\eqref{der:unrestricted-tangent}; applying $\RHom(S',-)$ to
\eqref{der:quot-sequence} gives the second identification.
In particular, forgetting the arrow and remembering its source has
differential the boundary
\begin{equation}\label{der:source-boundary}
\RHom(S',Q)\longrightarrow\RHom(S',S')[1].
\end{equation}

These are identifications of perfect complexes in families.  For
completeness, perfection and functoriality can also be checked before
fixing a vertex by taking pointed loops.  The relative loop stack at a
perfect family is its self-equivalence stack, open in the linear mapping
stack of its perfect endomorphism complex.  Its cotangent at the identity
is the dual endomorphism complex.  The cotangent triangle for the loop
fiber product identifies this with the pulled-back moduli cotangent
shifted by $[1]$.  Forgetting vertices and taking direct sums act on these
loop tangents by projection and block-diagonal sum.  Taking the fixed-target
fiber preserves these identifications.  This proves the claimed
functoriality, including the specified boundary.

Finally, at a geometric point $S'$ is torsion-free and $Q$ is torsion.
Surface Serre duality gives
\[
\Ext^2(S',Q)^\vee=\Hom(Q,S'(K))=0.
\]
Negative Ext groups vanish for sheaves; groups above degree two vanish
on a smooth surface.  Consequently the derived tangent has geometric
fibers in degrees $[0,1]$.  For a perfect complex the geometric-fiber
criterion is precisely the criterion for Tor-amplitude in this range.
\end{proof}

We will use the following comparison in both the absolute and relative
constructions.  We include the infinitesimal argument to specify which
cotangent morphism supplies the obstruction theory.

\begin{lemma}[Truncation comparison]\label{der:truncation}
Let $\mathfrak Z$ be a locally geometric derived stack, locally of finite
presentation, whose classical truncation $j:Z\to\mathfrak Z$ is a scheme.
Suppose $E_Z=j^*L_{\mathfrak Z}$ is perfect of Tor-amplitude $[-1,0]$.
The canonical map
\[
\phi_Z:E_Z\longrightarrow L_Z
\]
has cone in $D^{\leq-2}(Z)$, and is a perfect obstruction theory.
For a map of such enhancements, the induced cotangent map commutes with
the corresponding maps $\phi_Z$.
\end{lemma}

\begin{proof}
The cotangent map is induced by $j$, so its functoriality is that of
cotangent complexes.  The comparison statement is
\cite[Proposition~1.2 and Corollary~1.4]{STV}; the following argument gives
it directly in the present scheme-truncation situation.
Work on an affine open $x:\operatorname{Spec}R\to Z$, and write
$E=x^*E_Z$.  For every ordinary $R$-module $N$, split square-zero
extensions and agreement of $Z$ and $\mathfrak Z$ on classical rings give
an isomorphism
\[
\Hom(L_R,N)\xrightarrow{\sim}\Hom(E,N).
\]
There is also an injection
\[
\Hom(L_R,N[1])\longrightarrow\Hom(E,N[1]).
\]
Indeed, let a derivation class $\delta:R\to R\oplus N[1]$ map to zero
on the right.  The two points of $\mathfrak Z(R\oplus N[1])$ obtained
from $x$ by $\delta$ and by the zero derivation are then homotopic, over
$x$.  Infinitesimal cartesianness of a locally geometric derived stack
produces a lift of $x$ over the homotopy pullback
\[
R'=R\mathbin{\times^h_{R\oplus N[1]}}R,
\]
whose two maps are $\delta$ and zero.  This is the classical square-zero
extension of $R$ by $N$ classified by $\delta$.  The lift is a classical
point of $Z$ and factors through the same open $\operatorname{Spec}R$;
it therefore gives a splitting of $R'\to R$.  The extension class
$\delta$ is zero, proving injectivity.
Put $C=\Cone(E\to L_R)$.  The two tests, together with connectivity,
give $\Hom(C,N)=\Hom(C,N[1])=0$ for every $N$.  Since $C\in D^{\leq0}$,
the first equality implies $h^0(C)=0$, and the second then implies
$h^{-1}(C)=0$.  Thus $C\in D^{\leq-2}$.  Together with the
given perfect amplitude, this is the definition of a perfect obstruction
theory in \cite[Definitions~4.4 and~5.1]{BF}.
\end{proof}

\subsection{The based determinant and its amplitude}

The unrestricted Quot enhancement is now available. To impose the
determinant condition without introducing scalar automorphisms, we
first normalize line bundles at one point and then take a derived
fiber. The remaining amplitude check concerns its obstruction map.

Choose $o\in X$.  Let $\mathfrak P^o$ be the derived Picard stack of line
bundles on $X$ trivialized at $o$: it is the fiber, over the trivial line,
of evaluation at $o$ from the derived line-bundle stack to
$B\mathbb G_m$.  Its truncation is the Picard scheme $\Pic(X)$.
Indeed, every line $M$ on $X\times Y$ has the normalized form
\[
\nu(M)=M\otimes\pi^*(M|_{\{o\}\times Y})^{-1},
\]
with its canonical trivialization at $o$.  Two lines with the same
Picard class differ by a line from the base, which disappears on
normalization.  Moreover $\pi_*\mathcal O_{X\times Y}=\mathcal O_Y$,
so every line automorphism is a unit from the base; one preserving the
trivialization is the identity.  These observations identify the
truncation as a functor, not only on complex points.

The evaluation tangent is $\RGamma(\mathcal O_X)[1]\to\mathbb C[1]$.
Constant functions split it, and hence
\begin{equation}\label{der:based-picard}
T_{\mathfrak P^o}=\RGamma(\mathcal I_o)[1],\qquad
\RGamma(\mathcal O_X)=\mathbb C\oplus\RGamma(\mathcal I_o).
\end{equation}
In particular its cohomology lies in degrees zero and one, with groups
$H^1(\mathcal O_X)$ and $H^2(\mathcal O_X)$ respectively.

The derived determinant morphism of \cite[Section~3.1,
Proposition~3.2]{STV}, followed by $\nu$, gives a morphism
$\mathfrak U\to\mathfrak P^o$.  Set
\begin{equation}\label{der:derived-fiber}
\mathfrak B=
\mathfrak U\mathbin{\times^h_{\mathfrak P^o}}
\{\nu(L_1+L_2-D)\}.
\end{equation}
Its truncation is $B$: the fiber fixes the Picard class of $\det S'$
and hence of $\det Q$.  It makes no choice of trivialization of the
unnormalized determinant of the family at $o$.

The determinant differential is trace at geometric points
\cite[Proposition~3.2]{STV}. We need this identification on the
universal complexes over classical parameter schemes, including
nonreduced ones, and justify that family statement here.
Let $Y$ be such a quotient parameter scheme. The perfect source on
$X\times Y$ classifies a map to the derived stack of perfect complexes.
Pull back the differential of its determinant morphism to the
line-bundle stack $B\mathbb G_m$. The perfect-object tangent formula
gives an internal map
\[
 R\mathcal{H}om(S',S')[1]\longrightarrow\mathcal O_{X\times Y}[1].
\]
The determinant on complexes on $X$ is obtained by applying this
morphism to families on $X$; its differential is the derived pushforward
of this internal map.

Perfect trace duality identifies the internal map with an endomorphism
$a$ of the sheaf $S'$. On the locus where $S'$ is locally free, the
ordinary determinant differential gives $a=\id_{S'}$.
To extend this equality over $Y$, let $Z$ be the closed non-locally-free
locus of $S'$. Base flatness and the fiber criterion for local freeness
show that $Z$ has finite fibers over $Y$. Near a point of $Z$, choose
$f$ in its ideal whose image in the regular local surface fiber is a
nonzero element. That image is a nonzerodivisor; the local flatness
criterion makes $f$ a nonzerodivisor on the total space as well.
The image of $a-\id_{S'}$, viewed in $V_Y$, is coherent and supported
on $Z$, so locally a power of $f$ annihilates it. Since $V_Y$ is locally
free, this image is zero. Thus $a=\id_{S'}$ everywhere, proving the
trace identity as a map of complexes over $Y$. This argument applies
to the universal noetherian parameter schemes used below; naturality
and derived pullback preserve the identity on their base changes.

The tangent map defining \eqref{der:derived-fiber} is consequently the
composition of \eqref{der:source-boundary}, trace, and the projection
in \eqref{der:based-picard}.  We denote this composition by $d\det^o$.
The two complexes in this tangent fiber have amplitude $[0,1]$.
Their long exact sequence shows that its only possible extra group is
\[
 H^2(T_{\mathfrak B}|_b)=
 \coker\!\left\{\Ext^1(S',Q)\longrightarrow H^2(\mathcal O_X)\right\}
\]
at a geometric point $b$. The following rank bound makes this group
vanish.

\begin{lemma}[Surjectivity of the determinant obstruction map]
\label{der:trace-surjectivity}
Let $0\to F\to V_0\to Q_0\to0$ be an exact sequence on $X$, where
$V_0$ is a vector bundle of rank $r$ and $Q_0$ is torsion.  The map
\[
\Ext^1(F,Q_0)\longrightarrow\Ext^2(F,F)
\xrightarrow{\operatorname{tr}}H^2(\mathcal O_X)
\]
is surjective if $p_g=0$.  It is also surjective if $K$ is ample and
$Kc_1(Q_0)>rK^2$.
\end{lemma}

\begin{proof}
There is nothing to prove when $H^2(\mathcal O_X)=0$.  Otherwise Serre
duality and the exact sequence show that a section
$\omega\in H^0(K)$ annihilates the image precisely when
$\omega\id_F\in\Hom(F,F(K))$ lifts to a map
\[
a:V_0\longrightarrow F(K),\qquad a|_F=\omega\id_F.
\]
Writing $i:F\to V_0$, the maps $i(K)a$ and $\omega\id_{V_0}$ agree
where $F=V_0$, a dense open because $Q_0$ has rank zero.  They are maps
between vector bundles, so they agree everywhere.  Composing with the
quotient to $Q_0(K)$ gives $\omega Q_0=0$.

Suppose now that $\omega\ne0$, and write its effective divisor as
$C=(\omega)$.  The quotient $Q_0$ factors through $V_0|_C$.  At the
generic point of any integral curve $C_j\subset X$ the local ring is a
DVR.  If $\omega$ has order $a_j$ there, a quotient of
$(\mathcal O_{X,C_j}/(\omega))^r$ has length at most $ra_j$.
The divisorial cycle of a torsion sheaf is the sum of these generic
lengths, and represents its first Chern class: the elementary-divisor
decomposition over a DVR identifies this length with the order of the
determinant.  Therefore
\[
[Q_0]_{\mathrm{div}}\leq rC
\quad\hbox{componentwise},\qquad
Kc_1(Q_0)\leq rK^2.
\]
The latter inequality uses ampleness of $K$.  It contradicts the
hypothesis.  The argument allows embedded points and nonreduced curve
support, since it uses only generic lengths.  Thus no nonzero
$\omega$ annihilates the image, proving surjectivity.
\end{proof}

\begin{proposition}[The determinant-fixed Quot theory]\label{prop:quot-pot}
In either setting~\textup{(i)} or~\textup{(ii)}, the enhancement
$\mathfrak B$ induces a perfect obstruction theory
$E_B\to L_B$, with virtual tangent
\begin{equation}\label{eq:quot-tangent}
E_B^\vee\simeq
\fib\left\{
\RHom(S',Q)\xrightarrow{d\det^o}
\RGamma(\mathcal I_o)[1]\otimes\mathcal O_B
\right\}.
\end{equation}
It is $G$-equivariant.  Its virtual dimension is
\begin{equation}\label{eq:vdim}
v=2n+(A-K)D.
\end{equation}
\end{proposition}

\begin{proof}
The cotangent triangle of the derived fiber gives
\eqref{eq:quot-tangent} as an identity of perfect complexes.
Both its input and target have Tor-amplitude $[0,1]$; the input has
this amplitude by Lemma~\ref{der:arrow}.  Their fiber has no negative
cohomology and has no cohomology above degree two.  Its degree-two
cohomology vanishes exactly when
\[
\Ext^1(S',Q)\longrightarrow H^2(\mathcal O_X)
\]
is surjective at every geometric point.  In setting~(i) its target is
zero.  In setting~(ii), Lemma~\ref{der:trace-surjectivity} applies with
$r=2$, since $KD=3K^2>2K^2$.  The geometric-fiber criterion for perfect
complexes now gives amplitude $[0,1]$.  Lemma~\ref{der:truncation}
therefore supplies the specified obstruction morphism.

Postcomposition with the given $G$-action on $V$ defines an action on
the arrow stack, leaving the source determinant map invariant.  It
thus acts on the fiber \eqref{der:derived-fiber}, and naturality of the
cotangent map makes its obstruction theory equivariant.  Equivalently,
one can construct the truncation map on the quotient stacks relative
to $BG$, obtaining exactly this equivariant complex and morphism.

The Euler characteristic of the determinant target is $q-p_g=0$.
If $q_2=\int_X\ch_2(Q)$, Riemann--Roch gives
\[
\chi(V,Q)=2q_2-(K+L_1+L_2)D,\qquad
\chi(Q,Q)=-D^2.
\]
For the second identity, the only degree-four term in
$\ch(Q^\vee)\ch(Q)\td(X)$ is $-D^2$, since $Q$ has rank zero;
here $Q^\vee$ is the derived dual.
Consequently
\begin{align*}
v&=\chi(S',Q)-(q-p_g)
  =\chi(V,Q)-\chi(Q,Q)\\
 &=2\left(n+L_1D-\frac{D^2}{2}\right)
   -(K+L_1+L_2)D+D^2
  =2n+(A-K)D.
\end{align*}
\end{proof}

\subsection{Equivariant resolutions and the fixed theory}

The determinant-fixed obstruction theory is constructed. To apply
virtual localization, we still need global equivariant presentations
and a comparison identifying the fixed obstruction morphism itself.

\begin{lemma}[Equivariant global presentations]\label{der:global-resolution}
The scheme $B$ has a $G$-equivariant closed embedding in a smooth
projective scheme.  Its obstruction theory has a global presentation
by two $G$-equivariant vector bundles.  The morphism to the truncated
cotangent complex can be represented by a morphism of two-term
complexes for this embedding.
\end{lemma}

\begin{proof}
The Grassmannian construction of the projective Quot scheme is
equivariant for the action on $V$.  The determinant fiber is closed and
invariant, so restriction and a Pl\"ucker embedding give an equivariant
embedding $B\hookrightarrow M$ with $M$ smooth projective.  In particular
$B$ has a linearized ample bundle $\mathcal O_B(1)$.

Here is the global resolution argument for $E=E_B$.  For $m$ sufficiently
large, $h^0(E)(m)$ is generated by sections and
$H^a(h^{-1}(E)(m))=0$ for every $a>0$.  The hypercohomology sequence then lifts its
sections to $\mathbb H^0(E(m))$.  These are finite-dimensional torus
representations; linear reductivity allows an equivariant choice of
lifts.  Evaluation gives an equivariant bundle
$V_0=\mathcal O_B(-m)\otimes W_0$ and a map $V_0\to E$ surjective on
$h^0$.  Its derived fiber $V_1$ is perfect.  At every geometric point
$b$, the same map is surjective on $h^0$, since
$h^0(E\otimes^L k(b))=h^0(E)\otimes k(b)$.  Its fiber has cohomology
only in degree zero.  Thus $V_1$ is a vector bundle, and
$E\simeq[V_1\to V_0]$ in degrees $[-1,0]$.

Let $J$ be the conormal sheaf of $B\hookrightarrow M$, and set
$\Omega=\Omega_M|_B$.  The truncated cotangent complex
$\tau_{\geq-1}L_B$ is represented by $C=[J\to\Omega]$ in degrees
$[-1,0]$.  Increase $m$ above so that
$\Ext^1(V_0,J)=0$ as well; this is Serre vanishing for $J(m)$.
The triangle $J\to\Omega\to C\to J[1]$ then lifts
$V_0\to E\to C$ to $V_0\to\Omega$, equivariantly.  The map
$E\to C$ includes a nullhomotopy on its restriction from $V_1$.
Together with the chosen lift, that nullhomotopy gives a map
$V_1\to J=\fib(\Omega\to C)$.  These are ordinary sheaf maps and give
a commuting square
\[
\begin{array}{ccc}
V_1&\longrightarrow&V_0\\
\downarrow&&\downarrow\\
J&\longrightarrow&\Omega
\end{array}
\]
representing the original map to $C$.  All choices can be made in the
equivariant category, since taking torus invariants is exact.
\end{proof}

On the classical fixed locus, an invariant subsheaf of
$L_1\oplus L_2$ decomposes as the direct sum of its two character
subsheaves.  This holds over every classical base with trivial
$G$-action: the coaction decomposes a module into its weight modules,
and the inclusion in $V$ admits only the two indicated weights.
The kernels and quotients in these summands are flat because they are
direct summands of flat families.  Thus a fixed arrow is precisely a pair
\[
S'_1\longrightarrow L_1,\qquad S'_2\longrightarrow L_2,
\]
of rank-one kernel arrows with torsion quotients $Q_1,Q_2$, subject to
the total data and determinant constraint.  There are finitely many
discrete summand Hilbert polynomials on the fixed scheme: they are
locally constant, and the fixed scheme is of finite type.

Write $\mathfrak U_i$ for the unrestricted derived arrow Quot of
Lemma~\ref{der:arrow} with target $L_i$.  For each choice of the discrete
data, let
\[
\mathfrak W=(\mathfrak U_1\times\mathfrak U_2)
\mathbin{\times^h_{\mathfrak P^o}}
\{\nu(L_1+L_2-D)\},
\]
where the map to $\mathfrak P^o$ takes the normalized determinant of
the direct sum of the sources.  Its truncation $W$ is the corresponding
part of $B^G$, by the family decomposition just proved.

\begin{proposition}[The actual fixed obstruction theory]\label{prop:fixed-pot}
The enhancement $\mathfrak W$ induces a perfect obstruction theory
$E_W\to L_W$.  For $k:W\hookrightarrow B$, there is a natural
isomorphism of obstruction theories
\begin{equation}\label{der:fixed-comparison}
(Lk^*E_B)^{\mathrm{fix}}\xrightarrow{\sim}E_W,
\end{equation}
compatible with their maps to $L_W$.  Its virtual tangent is
\begin{equation}\label{der:fixed-tangent}
E_W^\vee\simeq
\fib\left\{
\bigoplus_{i=1}^2\RHom(S'_i,Q_i)
\longrightarrow\RGamma(\mathcal I_o)[1]\otimes\mathcal O_W
\right\}.
\end{equation}
The moving virtual normal complex is
\begin{equation}\label{der:moving-normal}
N_W^{\vir}=\bigoplus_{i\ne j}\RHom(S'_i,Q_j),
\end{equation}
where the $(i,j)$ summand has the difference between the characters on
$L_j$ and $L_i$, or evaluated weight $w_j-w_i$.
\end{proposition}

\begin{proof}
Direct sum is an actual morphism of perfect-arrow moduli problems.
The determinant of a direct sum is the tensor product of determinants,
compatibly with normalization.  Consequently direct sum induces a
derived map $\mathfrak k:\mathfrak W\to\mathfrak B$ lifting $k$.
This map is coherently equivariant.  Over the classifying stack $BG$,
tensor the source and target of each summand arrow with the same
character line of weight $+1$ or $-1$.  Normalizing the determinant
cancels each base character line against its restriction at $o$.
These canonical cancellations respect tensor products and the chosen
homotopy in the determinant fiber.  Thus direct sum defines
$BG\times\mathfrak W\to[\mathfrak B/G]$ over $BG$, with trivial
$G$-action on $\mathfrak W$.  Relative cotangents over $BG$ give its
equivariant cotangent morphism.
Its cotangent morphism, restricted to truncations, gives
\[
Lk^*E_B\longrightarrow E_W\longrightarrow L_W.
\]
By functoriality of the truncation comparison, this composite equals
$Lk^*E_B\to Lk^*L_B\to L_W$.

To identify this map, use the tangent formula for arrows.  On $W$,
\[
\RHom(S',Q)=\bigoplus_{i,j}\RHom(S'_i,Q_j).
\]
Direct sum maps the two input tangents to the diagonal blocks.  Trace
vanishes on off-diagonal blocks, and the based Picard tangent has weight
zero.  To display the comparison, abbreviate these complexes by
$\mathcal A_0=\bigoplus_i\RHom(S'_i,Q_i)$,
$\mathcal A=\bigoplus_{i,j}\RHom(S'_i,Q_j)$, and
$\mathcal P=\RGamma(\mathcal I_o)[1]\otimes\mathcal O_W$.
The derivative gives the morphism of fiber triangles
\[
\begin{array}{ccccc}
E_W^\vee&\longrightarrow&\mathcal A_0&\longrightarrow&\mathcal P\\
\downarrow&&\downarrow&&\Vert\\
Lk^*E_B^\vee&\longrightarrow&\mathcal A&\longrightarrow&\mathcal P,
\end{array}
\]
whose middle map is diagonal inclusion.  After taking weight zero, the
middle and right maps are identities.  The left map is therefore an
isomorphism of perfect complexes in
families.  The remaining summands are exactly
\eqref{der:moving-normal}.

In particular \eqref{der:fixed-tangent} has Tor-amplitude $[0,1]$,
as a direct summand of $Lk^*E_B^\vee$.  Lemma~\ref{der:truncation}
makes $E_W\to L_W$ a perfect obstruction theory.  Dualizing the
fiber-triangle comparison gives \eqref{der:fixed-comparison}.  The
compatibility already established shows that this is the fixed
obstruction morphism used in \cite[Proposition~1]{GP}, and not just an
identity of virtual tangent classes.
\end{proof}

The hypotheses of virtual localization are now in place: $B$ is
projective, Lemma~\ref{der:global-resolution} supplies an equivariant
embedding and global bundle presentations, and
Proposition~\ref{prop:fixed-pot} identifies the fixed obstruction
morphism.  Hence \cite[Section~3]{GP} gives
\begin{equation}\label{der:localization}
[B]^{\vir}=\sum_W k_*\frac{[W]^{\vir}}{\eu_G(N_W^{\vir})}
\quad\hbox{in }\CH_*^G(B)\otimes_{\mathbb Q[t]}\mathbb Q(t).
\end{equation}
Here $[W]^{\vir}$ is the ordinary virtual class just constructed, given
its trivial equivariant lift.  To see this explicitly, $G$ acts trivially
on $W$, and every equivariant coherent sheaf there decomposes into its
weights.  A weight-zero perfect theory is therefore the ordinary theory
with trivial weights.  On the mixed quotients defining equivariant Chow
groups, its intrinsic cone and zero-section Gysin construction are the
ordinary ones pulled back from $W$.  This identifies its equivariant
virtual class with that lift.  Every weight in the normal complex is
nonzero, so the Euler denominator is invertible after localization.
